%! Author = Admin
%! Date = 17.04.2026
\section{Описание математической модели объекта управления}

В работе в качестве объекта управления рассматривается квадрокоптер, движущийся в трехмерном пространстве вдоль заданной гладкой кривой.
Базовая непрерывная модель берется из диссертации Кима С.А. и реализована в модуле \texttt{code/drone\_sim/models/dynamics.py}. Вектор состояния расширенной модели имеет вид
\begin{gather*}
    x =
    \begin{bmatrix}
        p_x & p_y & p_z &
        v_x & v_y & v_z &
        \phi & \theta & \psi &
        \dot{\phi} & \dot{\theta} & \dot{\psi} &
        \bar{u}_1 & \rho_1 & u_2 & \rho_2
    \end{bmatrix}^{\top},
\end{gather*}
где $p = [p_x, p_y, p_z]^{\top}$ --- положение центра масс, $v = [v_x, v_y, v_z]^{\top}$ --- линейная скорость, $\phi,\theta,\psi$ --- углы рыскания, тангажа и крена соответственно. Переменные $\bar{u}_1,\rho_1,u_2,\rho_2$ описывают двойные интеграторы, вводимые при приведении модели к нормальной форме.

В принятой нормировке динамика квадрокоптера записывается как
\begin{gather*}
    \dot{p} = v,
\end{gather*}
\begin{gather*}
    \dot{v} = \frac{b(\phi,\theta,\psi)\bigl(u_1 + g\bigr) - [0,0,g]^{\top}}{m},
\end{gather*}
\begin{gather*}
    \ddot{\phi} = \frac{u_2}{J_{\phi}}, \qquad
    \ddot{\theta} = \frac{u_3}{J_{\theta}}, \qquad
    \ddot{\psi} = \frac{u_4}{J_{\psi}},
\end{gather*}
где $g$ --- ускорение свободного падения, $m$ --- масса, $J_{\phi},J_{\theta},J_{\psi}$ --- приведенные моменты инерции. Направление тяги определяется вектором
\begin{gather*}
    b(\phi,\theta,\psi) =
    \begin{bmatrix}
        \cos\phi \sin\theta \cos\psi + \sin\phi \sin\psi \\
        \sin\phi \sin\theta \cos\psi - \cos\phi \sin\psi \\
        \cos\theta \cos\psi
    \end{bmatrix}.
\end{gather*}

Для ограничения тяги используется гладкая функция насыщения
\begin{gather*}
    u_1 = \mathrm{sat}_L(\bar{u}_1) = L \tanh\left(\frac{\bar{u}_1}{L}\right),
\end{gather*}
что соответствует реализованной в коде функции \texttt{sat\_tanh}. Тогда уравнения для интеграторных расширений принимают вид
\begin{gather*}
    \dot{\bar{u}}_1 = \rho_1, \qquad \dot{\rho}_1 = v_1,
\end{gather*}
\begin{gather*}
    \dot{u}_2 = \rho_2, \qquad \dot{\rho}_2 = v_2,
\end{gather*}
а полный вектор управления равен
\begin{gather*}
    U = [v_1, v_2, u_3, u_4]^{\top}.
\end{gather*}

\subsection{Геометрия опорной траектории}

Заданная траектория описывается параметрически как гладкая кривая
\begin{gather*}
    p^{\ast}(s) = \begin{bmatrix} x^{\ast}(s) & y^{\ast}(s) & z^{\ast}(s) \end{bmatrix}^{\top},
\end{gather*}
где $s$ --- параметр кривой. В проекте используются кривые с постоянной нормой касательного вектора, так как именно для этого класса корректно работает алгоритм главы~4 диссертации:
\begin{gather*}
    \|t(s)\| = \left\|\frac{d p^{\ast}(s)}{d s}\right\| = \mathrm{const}.
\end{gather*}

В качестве типовых траекторий рассматриваются:
\begin{gather*}
    p^{\ast}_{\text{line}}(s) = \begin{bmatrix} s & s & s \end{bmatrix}^{\top},
\end{gather*}
\begin{gather*}
    p^{\ast}_{\text{circle}}(s) =
    \begin{bmatrix}
        r \cos \frac{s}{r} \\
        r \sin \frac{s}{r} \\
        0
    \end{bmatrix},
\end{gather*}
\begin{gather*}
    p^{\ast}_{\text{spiral}}(s) =
    \begin{bmatrix}
        r \cos s \\
        r \sin s \\
        k s
    \end{bmatrix}.
\end{gather*}

По кривой вычисляются геометрические характеристики, используемые в регуляторе и в ML-пайплайне:
\begin{gather*}
    \alpha(s) = \mathrm{atan2}\bigl(t_y(s), t_x(s)\bigr),
\end{gather*}
\begin{gather*}
    \beta(s) = \mathrm{atan2}\left(t_z(s), \sqrt{t_x^2(s) + t_y^2(s)}\right),
\end{gather*}
\begin{gather*}
    \varepsilon(s) = \frac{d \alpha(s)}{d s}.
\end{gather*}
Здесь $\alpha(s)$ задает угол рыскания касательной, $\beta(s)$ --- угол наклона касательной к горизонтальной плоскости, а $\varepsilon(s)$ характеризует кривизну проекции траектории на плоскость $XY$.

\subsection{Выходы системы и регулируемые ошибки}

Для траекторного движения по выходу вводится ближайшая точка на кривой с параметром $\zeta$. На ее основе формируется вектор регулируемых переменных
\begin{gather*}
    \lambda_1 =
    \begin{bmatrix}
        s_{\mathrm{arc}} - V^{\ast} t \\
        e_1 \\
        e_2 \\
        \delta_{\phi}
    \end{bmatrix},
\end{gather*}
где $V^{\ast}$ --- желаемая параметрическая скорость движения по кривой, $e_1,e_2$ --- ошибки в сопутствующей системе координат Френе, а
\begin{gather*}
    \delta_{\phi} = \phi - \phi^{\ast}(\zeta) = \phi - \alpha(\zeta)
\end{gather*}
--- ошибка по курсу. В актуальной реализации длина пройденной дуги вычисляется интегрально:
\begin{gather*}
    s_{\mathrm{arc}}(t) = \int\limits_{0}^{\zeta(t)} \|t(\tau)\|\,d\tau,
\end{gather*}
что корректно как для равномерной, так и для неравномерной параметризации кривой.

Наблюдатель ближайшей точки основан на условии ортогональности вектора ошибки касательной:
\begin{gather*}
    H(\zeta,p) = \bigl(p^{\ast}(\zeta) - p\bigr)^{\top} t(\zeta) = 0.
\end{gather*}
В численной схеме используется динамика
\begin{gather*}
    \dot{\zeta} = -\gamma_{\zeta} \,\rho \, H(\zeta,p),
\end{gather*}
где $\rho = \mathrm{sign}\left(\frac{dH}{d\zeta}\right)$, а коэффициент $\gamma_{\zeta}$ выбирается из условия устойчивости дискретного наблюдателя
\begin{gather*}
    0 < \gamma_{\zeta} \, \Delta t \, \|t\|_{\max}^{2} < 2.
\end{gather*}

\subsection{Ограничения модели}

Из кода и сопровождающей документации следуют основные практические ограничения модели:
\begin{enumerate}
    \item контроллер главы~4 гарантированно работает на кривых с постоянной нормой касательного вектора;
    \item при увеличении коэффициента усиления наблюдателя $\kappa$ необходимо уменьшать шаг интегрирования $\Delta t$;
    \item модель предполагает старт объекта в окрестности опорной кривой, иначе оценка ближайшей точки может сходиться медленно или уходить в неверную область.
\end{enumerate}

Таким образом, математическая модель задачи включает нелинейную динамику квадрокоптера, геометрию пространственной кривой, сопутствующую систему координат и наблюдатель параметра ближайшей точки. Именно в этой модели далее строится регулятор слежения по выходу и интеллектуальный модуль подбора скорости.
